\documentclass[12pt,a4paper]{article}

% ── Packages ──────────────────────────────────────────────
\usepackage[utf8]{inputenc}
\usepackage[T1]{fontenc}
\usepackage{amsmath,amssymb}
\usepackage{graphicx}
\usepackage{booktabs}
\usepackage[hidelinks]{hyperref}
\usepackage[margin=2.5cm]{geometry}
\usepackage{caption}
\usepackage{subcaption}
\usepackage{algorithm}
\usepackage{algpseudocode}
\usepackage{tikz}
\usepackage[numbers,sort&compress]{natbib}

% ── Title ─────────────────────────────────────────────────
\title{Reverse Engineering SaccadeJEPA:\\
  From Joint-Embedding Predictive Architectures\\
  to Gaze Dynamics Simulation}

\author{Emad Bahreinipour }
\date{May 2026}


% ══════════════════════════════════════════════════════════
\begin{document}
\maketitle

\tableofcontents
\newpage

% ══════════════════════════════════════════════════════════
% SECTION 1 — General Description of JEPA
% ══════════════════════════════════════════════════════════
\section{General Description of JEPA}
\label{sec:jepa}

\subsection{The Basic Idea}
\label{subsec:jepa-motivation}

A Joint-Embedding Predictive Architecture (JEPA) is a self-supervised learning
framework built around one central idea: instead of reconstructing missing data
in pixel space, the model predicts missing information in a learned
representation space. In plain terms, it does not try to reproduce every
texture, colour fluctuation, or noise pattern in an image. It tries to predict
the abstract state of what is missing.

This is different from pixel-reconstruction methods such as autoencoders, which
must model all visible detail, including information that may be irrelevant for
understanding the scene. It is also different from contrastive methods, which
often depend on negative examples and hand-crafted augmentations. JEPA instead
creates a predictive task directly from the input: observe one part, represent
another part, and train the model to predict the second representation from the
first.

\subsection{Core Principle: Prediction in Representation Space}
\label{subsec:jepa-core}

Let $\mathbf{x}$ be the observed part of an input and $\mathbf{y}$ the part to
be predicted. A context encoder maps the observed part to a representation,
while a target encoder maps the target part to another representation:

\begin{align}
  s_x &= f_\theta(\mathbf{x}),  \label{eq:enc-ctx} \\
  s_y &= f_{\bar\theta}(\mathbf{y}),  \label{eq:enc-tgt}
\end{align}

The predictor then estimates the target representation from the context
representation:

\begin{equation}
  \hat{s}_y = g_\phi(s_x), \label{eq:predictor}
\end{equation}

and training minimises the distance between the predicted and actual target
representations:

\begin{equation}
  \mathcal{L} = d\!\left(\hat{s}_y,\; s_y\right), \label{eq:loss}
\end{equation}

where $d(\cdot,\cdot)$ is usually a squared-error or Huber-style distance. The
important point is that the loss is not computed between images. It is computed
between embeddings. The model therefore learns to predict the abstract state of
the target, not its exact pixels.

\subsection{The EMA Target Encoder}
\label{subsec:jepa-collapse}

A practical issue in joint-embedding methods is collapse: both encoders could
learn to output nearly constant vectors, producing a low loss without useful
representations. JEPA avoids this by making the two branches asymmetric. The
context encoder is trained by gradient descent, but the target encoder is a
slow-moving copy of it.

Concretely, the target encoder parameters $\bar\theta$ are updated by an
exponential moving average (EMA) of the context encoder parameters $\theta$:

\begin{equation}
  \bar\theta \;\leftarrow\; \tau\,\bar\theta \;+\; (1-\tau)\,\theta,
  \label{eq:ema}
\end{equation}

with $\tau$ close to one. In plain terms, the target encoder is not trained
directly by back-propagation. It follows the context encoder slowly. This gives
the predictor a stable target to chase while still allowing the target
representation to improve over training.

\subsection{I-JEPA as the Image Example}
\label{sec:ijepa}
\label{subsec:ijepa-overview}

I-JEPA is the main image-based example of this idea. It divides an image into
patch tokens, hides large contiguous target blocks, and gives the context
encoder only the visible patches. The target encoder computes representations
for the hidden blocks, and the predictor learns to infer those target
representations from the visible context.

The important point for this report is not every architectural detail of
I-JEPA, but the pattern it establishes:

\begin{center}
\texttt{visible image region}
$\rightarrow$
\texttt{predict representation of unseen image region}.
\end{center}

Because I-JEPA predicts representations rather than pixels, it encourages the
encoder to learn semantic and spatially meaningful features. This makes it a
useful conceptual bridge to gaze: a fixation also reveals only part of the
visual scene, and the visual system must infer what lies outside the current
high-acuity region.

\subsection{Why JEPA Fits Gaze Modelling}
\label{subsec:ijepa-gaze-link}
\label{subsec:jepa-hierarchy}

The gaze interpretation follows naturally from the JEPA formulation. If
$\mathbf{x}$ is the current fixation crop and $\mathbf{y}$ is the next fixation
crop, then the task becomes:

\begin{center}
\texttt{current fixation representation + movement information}
$\rightarrow$
\texttt{next fixation representation}.
\end{center}

This is the bridge used in the rest of the report. Original SaccadeJEPA already
applies JEPA to spatial shifts: it predicts the representation of a shifted
crop from another crop and an encoded movement. The adaptation in this report
keeps that prediction machinery, but changes where the shifts come from.
Instead of one random perturbation, the model receives saliency-driven gaze
transitions.
% ══════════════════════════════════════════════════════════
% SECTION 2 — Adaptation for attentive vision / gaze
% ══════════════════════════════════════════════════════════
\section{Adaptation for Attentive Vision and Gaze}
\label{sec:saccade}

\subsection{What Original SaccadeJEPA Provides}
\label{subsec:saccade-overview}

The starting point is the \texttt{SaccadeJepa} model from the LumenPallidium
JEPA repository~\cite{lumenpallidium2024jepa}. It is an unpublished JEPA-style
variant, separate from the repository's standard \texttt{ViTJepa}
implementation. Its pretext task is simple: given two spatially related crops
from the same image, predict the target encoder's representation of one crop
from the context encoder's representation of the other.

In plain words, original SaccadeJEPA does this:

\begin{center}
\texttt{image -> one random crop pair -> encode -> predict target representation}.
\end{center}

The main components are already useful for gaze modelling:

\begin{itemize}
  \item \textbf{Context encoder:} a ConvNeXt-Tiny network that encodes the
        first crop.
  \item \textbf{Target encoder:} an EMA copy of the context encoder that
        encodes the second crop without gradient updates.
  \item \textbf{Affine embedder:} a NeRF-style sinusoidal embedding of the
        spatial displacement between the two crops.
  \item \textbf{Predictor:} a three-layer MLP that receives the context
        representation plus the movement embedding and predicts the target
        representation.
\end{itemize}

The original cropper is the weak part for gaze. \texttt{SaccadeCropper}
samples one random affine transformation per image. The sampled movement is not
conditioned on image content, does not form a sequence, and does not remember
where the model has already looked. It is useful as a representation-learning
pretext task, but it is not yet an explicit gaze simulator.

\subsection{What Needs to Change for Gaze}
\label{subsec:saccade-reuse}

The biological interpretation of gaze requires three behaviours that original
SaccadeJEPA does not provide:

\begin{itemize}
  \item \textbf{Content-driven target selection.} The next crop should depend
        on the image, not on a uniform random draw.
  \item \textbf{Sequential fixations.} The model should produce a short
        scanpath, not only one independent crop pair.
  \item \textbf{Inhibition of return.} Recently visited locations should be
        temporarily suppressed so the model does not repeatedly sample the same
        point.
\end{itemize}

The conservative design choice is to change only the component responsible for
choosing and producing crops. The encoder, target encoder, EMA update,
predictor, affine embedder, and loss are kept from SaccadeJEPA. This keeps the
adaptation easy to interpret: if the patched model learns, the learning signal
comes from the same JEPA machinery, but applied to gaze-like transitions.

\subsection{The Cropper Substitution}
\label{subsec:loop:motivation}
\label{subsec:loop:method}

The adaptation is implemented in the project's JEPA-reuse module. Its factory
first builds a normal \texttt{SaccadeJepa} instance and then replaces one
attribute:

\begin{center}
\texttt{sj.saccade\_cropper = GazeCropper(...)}
\end{center}

This is the central architectural change of the project. Before the swap, the
model receives one random crop pair. After the swap, it receives a sequence of
saliency-driven fixation transitions:

\begin{center}
\texttt{image -> saliency map -> fixation 1 -> fixation 2 -> ... -> JEPA prediction}.
\end{center}

\texttt{GazeCropper} preserves the same return signature as the upstream
\texttt{SaccadeCropper}:

\begin{equation}
  (\mathbf{view}_1,\; \mathbf{view}_2,\; \mathbf{affines})
  = \texttt{cropper}(\mathbf{x}).
  \label{eq:gaze-cropper-return}
\end{equation}

Because the rest of \texttt{SaccadeJepa.forward} sees the same three tensors,
the upstream forward pass does not need to be rewritten.

\subsection{How GazeCropper Works}
\label{subsec:saccade-generation}

For each image, \texttt{GazeCropper} performs a short explicit gaze loop:

\begin{enumerate}
  \item Compute a saliency map from the image.
  \item Initialise an inhibition-of-return mask to ones.
  \item For $T$ steps, sample the next fixation from
        $\mathrm{saliency} \times \mathrm{IOR}$.
  \item After each fixation, suppress that location with a Gaussian inhibition
        well and let older suppressions decay.
  \item Convert the $T$ fixations into $T-1$ transitions
        $(\mathrm{fix}_t \to \mathrm{fix}_{t+1})$.
\end{enumerate}

In the experiments below, $T=5$, so each image produces four transitions. Each
transition contributes a $128 \times 128$ crop at the current fixation, a
$128 \times 128$ crop at the next fixation, and the normalised displacement
between them. The predictor is therefore trained on:

\begin{center}
\texttt{current fixation crop + movement embedding}
$\rightarrow$
\texttt{next fixation representation}.
\end{center}

The saliency source is pluggable: any module that maps an RGB image batch to a
normalised saliency map can be used. This includes random saliency, centre
bias, classical spectral-residual saliency, or Arash Khosropour's trained
ResNet saliency model. The ResNet model is the default for the final
experiments; it is trained separately on FIND fixation maps and then frozen
during JEPA training.

\subsection{Compatibility with SaccadeJEPA}
\label{subsec:saccade-forward}

The substitution is safe because \texttt{GazeCropper} respects the two
contracts expected by the upstream model.

\paragraph{First, the embedding dimension is unchanged.}
The upstream affine embedder expects a 36-dimensional NeRF-style movement
embedding. This comes from translation order $L_t=4$ and rotation order
$L_r=10$:

\begin{equation}
  \texttt{affine\_embed\_dim} = L_t \cdot 4 + L_r \cdot 2 = 36.
  \label{eq:nerf-encoding}
\end{equation}

\texttt{GazeCropper} reads these orders from the original cropper and produces
the same embedding width. Rotation is set to zero, because gaze shifts translate
the fixation window but do not rotate the world.

\paragraph{Second, the encoder inputs remain symmetric.}
Both the context encoder and the target encoder still receive $3$-channel
$128 \times 128$ RGB crops. This matters because the target encoder is an EMA
copy of the context encoder. An earlier integrated-controller design used a
$6$-channel fovea-plus-periphery context branch with a $3$-channel target
branch; that would have made the EMA path asymmetric. The final design avoids
that issue by keeping the upstream $3$-channel crop interface.

\subsection{Losses Reused from SaccadeJEPA}
\label{subsec:saccade-losses}

The training objective is also reused. The main term compares the predictor's
output with the EMA target encoder's representation of the next fixation crop:

\begin{equation}
  \mathcal{L}_{\text{pred}}
  = \texttt{HuberLoss}\!\bigl(\hat{s}_y,\; s_y\bigr).
  \label{eq:saccade-pred-loss}
\end{equation}

The upstream cycle-consistency term asks the predictor to approximately undo a
transition:

\begin{equation}
  \mathcal{L}_{\text{cycle}}
  = \bigl\lVert \hat{s}_x - s_x \bigr\rVert_2^2.
  \label{eq:saccade-cycle-loss}
\end{equation}

VICReg variance and covariance terms are kept as collapse-prevention
regularisers. In compact form, the total objective is:

\begin{equation}
  \mathcal{L}
  = \mathcal{L}_{\text{pred}}
  + \lambda_{\text{cycle}}\,\mathcal{L}_{\text{cycle}}
  + \lambda_{\text{var}}\,\mathcal{L}_{\text{var}}
  + \lambda_{\text{cov}}\,\mathcal{L}_{\text{cov}}.
  \label{eq:saccade-total-loss}
\end{equation}

Thus the adaptation is deliberately minimal. The source of the two views
changes from a random affine perturbation to a saliency-driven gaze transition;
the JEPA prediction machinery remains the same.

\begin{figure}[t]
  \centering
  \includegraphics[width=\linewidth]{../outputs/jepa_reuse_demo.png}
  \caption{The patched \texttt{SaccadeJepa} on a FIND frame at
    $200 \times 200$. \texttt{GazeCropper} consumes the frame and a
    saliency map, runs $T=5$ saliency~$\times$~IOR fixations, and emits four
    $(\mathbf{view}_1, \mathbf{view}_2, \mathbf{affine})$ transitions in the
    same return format expected by upstream SaccadeJEPA. Every box after the
    cropper is reused from the original model.}
  \label{fig:loop:in-action}
\end{figure}

% ══════════════════════════════════════════════════════════
% SECTION 3 — Possible simulations and results
% ══════════════════════════════════════════════════════════
\section{Possible Simulations and Results}
\label{sec:loop}

The feasible simulation is a self-supervised gaze-transition prediction task
on FIND~\cite{liu2016find}. The model receives a static video frame, uses
\texttt{GazeCropper} to generate a short saliency-driven scanpath, and trains
SaccadeJEPA's predictor to anticipate the representation of the next fixation
crop. This does not yet claim to reproduce human scanpaths exactly. It tests a
more basic question: after the cropper substitution, can the original JEPA
machinery still learn?

\subsection{Simulation Setup}
\label{subsec:loop:evaluation}

Frames are decoded from FIND raw MP4 videos and resized to
$200 \times 200$. The shared split is 50 training videos, 6 validation videos,
and 6 test videos. For the demonstration run, each training video contributes
four sampled frames per epoch. With batch size 4 and $T=5$ fixations per
frame, the predictor sees 16 fixation transitions per batch.

Training uses the same loss family described in
Sec.~\ref{subsec:saccade-losses}: Huber prediction loss~\cite{huber1964robust},
cycle consistency, and VICReg variance/covariance regularisation
~\cite{bardes2022vicreg}. AdamW is used with learning rate
$5 \cdot 10^{-5}$ and EMA decay $\tau=0.996$~\cite{grill2020byol}. The only
practical adjustment is that the VICReg weights are reduced for the small CPU
batch size: $\lambda_{\mathrm{var}}=1.0$,
$\lambda_{\mathrm{cov}}=0.1$, and $\gamma=1.0$. The upstream defaults were
tuned for a much larger effective batch.

The run lasts five CPU epochs. Figure~\ref{fig:loop:loss-curve} shows that
both total loss and prediction MSE decrease over training. The key curve is
prediction MSE, because it measures how close the predictor is to the EMA
target representation of the next fixation.

\begin{figure}[t]
  \centering
  \includegraphics[width=0.92\linewidth]{../outputs/jepa_reuse/metrics/loss_curve.png}
  \caption{Training loss and prediction MSE for the patched
    \texttt{SaccadeJepa} on FIND. Five epochs, batch size~$4$,
    four frames per training video, on CPU. The total loss combines
    a Huber prediction term, a cycle term, and the VICReg
    variance/covariance regulariser. The prediction MSE is read
    directly from the predictor against the EMA target encoder.
    Per-batch values are shown lightly; the heavy line is a running
    mean.}
  \label{fig:loop:loss-curve}
\end{figure}

\subsection{Quantitative Results}
\label{subsec:loop:results}

The main result is the validation prediction MSE. Before training, the patched
model gives a validation MSE of $0.2996$. After five epochs, the matched
validation pass gives $0.0462$. This is an $84.6\%$ reduction. The result means
that the predictor learned to anticipate the next-fixation representation under
the saliency-driven cropper.

\begin{table}[h]
  \centering
  \begin{tabular}{lr}
    \toprule
    Phase & Prediction MSE \\
    \midrule
    Untrained, validation pool ($n_{\mathrm{val}}{=}24$)        & $0.2996$ \\
    \addlinespace
    Validation, end of epoch~$1$                                & $0.0963$ \\
    Validation, end of epoch~$2$                                & $0.0471$ \\
    Validation, end of epoch~$3$                                & $0.0459$ \\
    Validation, end of epoch~$4$                                & $0.0428$ \\
    Validation, end of epoch~$5$                                & $0.0447$ \\
    \addlinespace
    Trained, validation pool (matched val pass on \texttt{final.pt}) & $0.0462$ \\
    \bottomrule
  \end{tabular}
  \caption{Predictor mean-squared error against the EMA target
    encoder, before and after a five-epoch training run on FIND.
    All numbers are averaged across the six held-out validation
    videos sampled at four frames per video (batch size~$4$,
    $T{=}5$ fixations per frame), using Arash's trained
    \texttt{ResNetSaliency} as the saliency source. The last row
    re-loads the final checkpoint and replays the validation pass
    in matched conditions: the training-time number recorded at
    end of epoch~$5$ ($0.0447$) and the post-hoc reproduction
    ($0.0462$) agree to within batch-level noise.}
  \label{tab:loop:ks}
\end{table}

\subsection{Qualitative Scanpath Simulation}

Figure~\ref{fig:loop:qualitative} shows four scanpaths sampled on held-out FIND
test frames. The cropper itself is not trained by the JEPA loss; its fixations
come from the frozen saliency source and the IOR mask. The figure is therefore
not evidence that the model has learned a human scanpath policy. Its role is
more limited: it shows that the substituted cropper produces structured,
sequential fixations on real video frames, and that the crops passed into
SaccadeJEPA are visually meaningful.

\begin{figure}[t]
  \centering
  \includegraphics[width=\linewidth]{../outputs/jepa_reuse_scanpaths_trained.png}
  \caption{Fixation sequences produced by the trained
    \texttt{GazeCropper} on four FIND test frames. Numbered red
    markers are the five fixations in temporal order; lines connect
    consecutive fixations. The saliency source is Arash's trained
    \texttt{ResNetSaliency} (ResNet-18 backbone with a
    $1{\times}1$-convolution head, drawn from the shared
    \texttt{gazejepa.saliency} package), and the inhibition mask
    uses the same
    $\sigma_{\text{IOR}}{=}20$~px and decay~$0.7$ as in the
    training run.}
  \label{fig:loop:qualitative}
\end{figure}

\subsection{Interpretation}
\label{subsec:loop:discussion}

These results answer the feasibility question positively. With a one-component
substitution, SaccadeJEPA can be made to operate on explicit gaze transitions,
and the original prediction machinery still converges. The decrease in MSE is
the same kind of JEPA signal as in the original model: the predictor approaches
the EMA target representation. The difference is that the two views now come
from consecutive saliency-driven fixations rather than one random affine crop.

The result should be interpreted carefully. The model is not trained with a
human scanpath loss, and it does not receive gradients through the discrete
fixation choices. It is therefore better described as a reusable JEPA backbone
for gaze-transition prediction, not as a complete learned gaze policy.

\subsection{Limitations}
\label{subsec:loop:limitations}

The main limitations are:

\begin{itemize}
  \item \textbf{Static frames.} FIND is a video dataset, but this simulation
        runs $T=5$ fixations on a single frozen frame.
  \item \textbf{No fixation duration.} The model predicts where the next crop
        is, not how long the eye stays there.
  \item \textbf{Frozen saliency.} Arash's \texttt{ResNetSaliency} is used as a
        fixed prior. It is not updated by the JEPA loss.
  \item \textbf{Small training budget.} The run is five epochs on CPU, so it is
        a feasibility check rather than a tuned final model.
  \item \textbf{No learned policy.} A future version could learn the fixation
        policy directly, for example with a recurrent attention model in the
        spirit of Mnih et al.~\cite{mnih2014ram}.
\end{itemize}

% ══════════════════════════════════════════════════════════
% SECTION 4 — Conclusion
% ══════════════════════════════════════════════════════════
\section{Conclusion}
\label{sec:conclusion}

This report asked whether SaccadeJEPA can be reused in a less obscure way for
explicit gaze simulation. The answer is yes, within a carefully limited scope.
Original SaccadeJEPA already contains the useful JEPA machinery: a context
encoder, an EMA target encoder, a movement embedding, a predictor, and a
representation-level prediction loss. Its limitation is the source of the
views: one random crop pair per image.

The proposed adaptation changes that source of views. Replacing
\texttt{SaccadeCropper} with \texttt{GazeCropper} turns the model from:

\begin{center}
\texttt{image -> random crop pair -> predict representation}
\end{center}

into:

\begin{center}
\texttt{image -> saliency-driven fixation sequence -> predict next fixation}.
\end{center}

The experiment on FIND shows that this substitution is compatible with the
upstream loss: validation prediction MSE drops from $0.2996$ to $0.0462$ after
five CPU epochs. This is not a claim that the model fully matches human gaze,
but it is strong feasibility evidence that SaccadeJEPA can be reused as a
representation-space predictor over explicit gaze transitions.


\begin{thebibliography}{99}

\bibitem{bardes2022vicreg}
A.~Bardes, J.~Ponce, and Y.~LeCun.
\newblock {VICReg}: {V}ariance-{I}nvariance-{C}ovariance regularization for
  self-supervised learning.
\newblock In \emph{International Conference on Learning Representations
  (ICLR)}, 2022.

\bibitem{grill2020byol}
J.-B. Grill, F.~Strub, F.~Altch\'e, C.~Tallec, P.~H. Richemond, E.~Buchatskaya,
  C.~Doersch, B.~Avila~Pires, Z.~D. Guo, M.~Gheshlaghi~Azar, B.~Piot,
  K.~Kavukcuoglu, R.~Munos, and M.~Valko.
\newblock Bootstrap your own latent: a new approach to self-supervised
  learning.
\newblock In \emph{Advances in Neural Information Processing Systems
  (NeurIPS)}, 2020.

\bibitem{hou2007spectralresidual}
X.~Hou and L.~Zhang.
\newblock Saliency detection: a spectral residual approach.
\newblock In \emph{Proceedings of the {IEEE} Conference on Computer Vision and
  Pattern Recognition (CVPR)}, 2007.

\bibitem{huber1964robust}
P.~J. Huber.
\newblock Robust estimation of a location parameter.
\newblock \emph{Annals of Mathematical Statistics}, 35(1):73--101, 1964.

\bibitem{itti1998saliency}
L.~Itti, C.~Koch, and E.~Niebur.
\newblock A model of saliency-based visual attention for rapid scene analysis.
\newblock \emph{IEEE Transactions on Pattern Analysis and Machine
  Intelligence}, 20(11):1254--1259, 1998.

\bibitem{klein2000ior}
R.~M. Klein.
\newblock Inhibition of return.
\newblock \emph{Trends in Cognitive Sciences}, 4(4):138--147, 2000.

\bibitem{liu2016find}
Y.~Liu and M.~Xu.
\newblock {FIND}: a database for multi-face tracking with fixation data on
  multiparty conversation videos.
\newblock In \emph{Proceedings of the {IEEE} International Conference on
  Multimedia and Expo (ICME)}, 2016.

\bibitem{liu2022convnext}
Z.~Liu, H.~Mao, C.-Y. Wu, C.~Feichtenhofer, T.~Darrell, and S.~Xie.
\newblock A {ConvNet} for the 2020s.
\newblock In \emph{Proceedings of the {IEEE/CVF} Conference on Computer Vision
  and Pattern Recognition (CVPR)}, 2022.

\bibitem{lumenpallidium2024jepa}
{LumenPallidium}.
\newblock {JEPA}: a {PyTorch} implementation of joint-embedding predictive
  architectures, including {SaccadeJEPA}.
\newblock Software repository, 2024.
\newblock \url{https://github.com/LumenPallidium/jepa}, accessed 2026.

\bibitem{mnih2014ram}
V.~Mnih, N.~Heess, A.~Graves, and K.~Kavukcuoglu.
\newblock Recurrent models of visual attention.
\newblock In \emph{Advances in Neural Information Processing Systems
  (NeurIPS)}, 2014.

\bibitem{posner1984ior}
M.~I. Posner and Y.~Cohen.
\newblock Components of visual orienting.
\newblock In H.~Bouma and D.~G. Bouwhuis, editors, \emph{Attention and
  Performance X: Control of Language Processes}, pages 531--556. Erlbaum, 1984.

\end{thebibliography}

\end{document}
